\section{Custom gates: packaging a circuit as one block}
\hypertarget{custom-gates}{}

A custom gate is a whole QPU process saved under a short name so that it can
be placed like any built-in block. It adds no new physics: when it runs, the
simulator replays the process's own gates, one by one, on the wires you
chose. Understanding a custom gate therefore means understanding the smaller
gates inside it.

\subsection{Register one in four steps}

\begin{enumerate}
  \item Paste the process below into the protocol editor in
  \textbf{Compile text into a circuit} and choose \textbf{Compile protocol}
  to confirm it has no errors.
  \item In \textbf{Register custom gates from QPU processes}, type
  \code{NOR} in \textbf{Gate id (palette label)} and leave
  \textbf{Source process} on \emph{Current protocol editor}.
  \item Choose \textbf{Register custom gate}. A \textbf{NOR} block appears
  under \textbf{Custom gates} in the palette and in the workbench
  \textbf{Gate} list.
  \item Place it with \textbf{Control A}, \textbf{Control B}, and
  \textbf{Target particle}, exactly like AND.
\end{enumerate}

\textbf{Complete runnable protocol}
\begin{lstlisting}
PARAMS: A:state B:state

MAIN-PROCESS NorGate
CREATETOKEN -I Out

OR -I A B -O Out
X -I Out -O Out

RETURNVALS Out
\end{lstlisting}

\subsection{How the process maps onto wires}

\begin{center}
\begin{tikzpicture}[font=\small]
  \foreach \y/\q in {1/q0,0/q1,-1/q2} {
    \draw[thick] (0,\y)--(3,\y);
    \node[left] at (0,\y) {\q};
    \draw[thick] (5,\y)--(10,\y);
    \node[left] at (5,\y) {\q};
  }
  \node[draw=QpuCyan,fill=QpuCyan!12,minimum width=1.1cm,minimum height=2.5cm,font=\bfseries] at (1.5,0) {NOR};
  \node[font=\scriptsize,anchor=south east] at (0.92,1) {A};
  \node[font=\scriptsize,anchor=south east] at (0.92,0) {B};
  \node[font=\scriptsize,anchor=south east] at (0.92,-1) {Out};
  \node at (4,0) {\Large$=$};
  \fill[QpuNavy] (6.5,1) circle (.09);
  \fill[QpuNavy] (6.5,0) circle (.09);
  \draw[thick] (6.5,1)--(6.5,-1);
  \node[draw=QpuBlue,fill=QpuBlue!8,minimum width=.72cm,minimum height=.5cm,font=\scriptsize\bfseries] at (6.5,-1) {OR};
  \node[draw=QpuBlue,fill=QpuBlue!8,minimum width=.6cm,minimum height=.5cm,font=\bfseries] at (8.5,-1) {X};
\end{tikzpicture}
\end{center}

\begin{center}
\begin{tabularx}{\textwidth}{>{\bfseries}p{0.27\textwidth}X}
\toprule
In the process & Becomes, when the block is placed\\
\midrule
1st \code{PARAMS} entry & \textbf{Control A} (the lowest free wire when
dropped on the canvas).\\
2nd \code{PARAMS} entry & \textbf{Control B}. Further entries take the next
free wires.\\
1st \code{RETURNVALS} entry & \textbf{Target particle}: the only wire the
block should change.\\
Any other token & A fresh extra workspace wire that the block allocates for
itself.\\
\bottomrule
\end{tabularx}
\end{center}

The palette gives every custom block exactly one target wire, so register
processes with a single \code{RETURNVALS} output. A process with two outputs,
such as a half-adder, compiles on its own but cannot be placed as a block; use
it through \code{RUNCHILD} instead.

\subsection{Reasoning about reversibility}

A sequence of reversible gates is reversible: undo the last gate first, then
the one before it, and so on.
\[
(U_k\cdots U_2U_1)^{-1}=U_1^{-1}U_2^{-1}\cdots U_k^{-1}.
\]
NorGate is OR followed by X. Both undo themselves, so NOR undoes itself too:
placing the block twice in a row returns every wire to where it started.
Written as a rule, the block performs
\[
\mathrm{Out}'=\mathrm{Out}\oplus\neg(A\lor B),
\]
the same ``target XOR function'' shape as AND. That shape is always
reversible, because applying it again XORs the same value back out.

\begin{cautionbox}[Three things that make a custom gate irreversible]
\begin{description}
  \item[MEASURE inside the process] The block collapses its wires every time
  it runs.
  \item[SET on the output] \code{SET Out 0p} makes the compiler insert a
  hidden RESET, so the block \emph{overwrites} the target
  ($\mathrm{Out}'=f(A,B)$) instead of XORing into it. Whatever the target held
  before is erased. NorGate deliberately has no SET: a freshly compiled wire
  already starts at 0, so it passes the same truth table.
  \item[Temporary wires left dirty] Helper tokens become hidden workspace
  wires. If a helper is not returned to 0 (uncomputed), it stays entangled
  with the inputs, and the visible wires alone no longer behave reversibly.
\end{description}
\end{cautionbox}

\begin{beginnerbox}[A quick test you can run in the builder]
Place your block twice on the same wires and choose \textbf{Run all}. For a
reversible out-of-place gate, every wire ends exactly as it started. Try it
with the target's \emph{name} start set to \code{0p} and then \code{1p}. If
the target does not return, look for one of the three problems above.
\end{beginnerbox}

\subsection{Rules to remember}
\begin{itemize}
  \item The gate id must start with a letter, use only letters, digits, and
  underscores, and cannot match a built-in gate such as \code{AND}.
  \item Custom gates are saved for the current browser session only, not in
  the bundled catalog. Keep the \code{.qpucir} file to recreate one later.
  \item A custom gate can wrap a process that itself uses \code{RUNCHILD};
  the child processes are captured when you register it.
\end{itemize}
